%chapter2.tex
\chapter{格式要求}
\label{geshiyaoqiu}
学位论文每部分一级标题从新的一页开始，各部分要求如下：

\section{封面}
封面是学位论文的外表面，对论文起装潢和保护作用，并提供相关的信息。博士学位论文、硕士学位论文、专业博士学位论文、专业硕士学位论文使用不同封面。

学校代码：10459

学号或申请号：全日制研究生填写学号，在职申请学位人员填写申请号。

密级：非涉密（公开）论文不需标注密级，涉密论文须标注论文的密级（内部、秘密、或机密），同时还应注明相应的保密年限。

论文题目：应简明扼要地概括和反映出论文的核心内容，一般不宜超过25字，必要时可加副标题，若本论文为国家自然基金等资助项目可在题目下一行用宋体12磅居中注明。

培养院系：填写所属院（系）的全名。

学科门类、专业名称：以国务院学位委员会1997年颁布的《授予博士、硕士学位和培养研究生的学科、专业目录》为准，学科门类以十二大门类为准填写，专业名称以二级学科为准填写。

专业学位名称（仅限于专业博士、硕士学位论文）：填写所申请的专业硕士学位的全称，如工商管理硕士、工程硕士等。

导师姓名：填写导师的姓名、职称（教授、研究员等）

完成时间：填写论文成文打印日期。格式：例2010年5月；2010年12月

\section{英文封面}
按照中文封面格式在封面相应位置填写英文名称，含中文封面的主要内容。

\section{学位论文原创性声明}
本部分内容直接从郑州大学研究生院资料下载中下载，提交时学位论文作者须亲笔签名。

\section{学位论文使用授权声明}
本部分内容直接从郑州大学研究生院资料下载中下载，提交时学位论文作者须亲笔签名。

\section{摘要}
应概括地反映出本论文的主要内容，具有独立性和自创性，包括研究工作目的、方法、结果和结论，要突出本论文的创造性成果。摘要力求语言精炼准确，硕士学位论文建议1000字以内，博士学位论文建议2000字以内。部分学生用外文撰写学位论文时，博士学位论文的摘要应不少于5000字符，硕士应不少于2500字符。摘要中不可出现图片、图表、表格或其他插图材料。

关键词是为了便于做文献索引和检索工作而从论文中选取出来用以表示全文主题内容信息的单词或术语。关键词应体现论文特色，具有语义性，在论文中有明确的出处，并应尽量采用《汉语主题词表》或各专业主题词表提供的规范词。在摘要内容下方另起一行标明，一般3～8个，之间用‘；’分开。为便于国际交流，应标注与中文对应的英文关键词。


\section{Abstract}
Abstract内容与中文摘要相对应，应注意外文大小写、正斜体书写符合有关要求。

\section{目录}
目录应将文内的章节标题依次排列。

\section{图和附表清单}
图表较多时使用。图的清单应有序号、图题和页码。表的清单应有序号、表题和页码。

\section{符号说明}
如果论文中使用了大量的符号、标志、缩略词、首字母缩写、专门计量单位、自定义名词和术语等，应编写成注释说明汇集表。若上述符号使用数量不多，可以不设此部分，但必须在论文中出现时加以说明。

\section{正文}
正文是学位论文的主体和核心部分，每一章一级标题应另起页，一般包括以下几个方面：

（1）引言（第1章）：包括研究的目的和意义，问题的提出，选题的背景，文献综述，研究方法，论文结构安排等。

（2）各具体章节：本部分是论文作者的研究内容，是论文的核心。各章之间互相关联，符合逻辑顺序。

（3）结论（最后1章）：是学位论文最终和总体的结论，不是正文中各段小结的简单重复，应明确、精练、完整、准确。着重阐述作者的创造性工作及所取得的研究成果在本学术领域的地位、作用和意义，结论应包括论文的核心观点，交代研究工作的局限，提出未来工作的意见或建议。


\section{注释}
当论文中的字、词或短语，需要进一步加以说明，而又没有具体的文献来源时，用注释。注释一般在社会科学中用得较多。应控制论文中的注释数量，不宜过多。由于论文篇幅较长，建议采用文中编号加“脚注”的方式。

\section{参考文献}
参考文献表是文中引用的有具体文字来源的文献集合。为了反映论文的科学依据和作者尊重他人研究成果的严肃态度以及向读者提供有关信息的出处，应列出参考文献表。参考文献表中列出的一般应限于作者直接阅读过的、最主要的、发表在正式出版物上的文献。私人通信和未公开发表的资料，一般不宜列入参考文献，可紧跟在引用的内容之后在文中编号，用“脚注”的方式标注在当页的下方。正文中引用参考文献的部位，须用上标标注[参考文献序号]，如 \cite{chen2001,nadkarni1992,hua1973,zhu1973,huo1981,timoshenko1959,zhang1991,dupont1974,ding2001,cairns1965,patent,standard,Karl2022}。参考文献应按文中引用出现的顺序排列。

\section{附录}
有些材料编入文章主体会有损于编排的条理性和逻辑性，或有碍于文章结构的紧凑和突出主题思想等，可将这些材料作为附录编排于全文的末尾。

附录的序号用A，B，C…系列，如附录A，附录B…。附录中的公式、图和表的编号分别用A1，A2…系列；图A1，图A2…系列；表A1，表A2…系列。每个附录应有标题。


附录的序号用A，B，C…系列，如附录系列，如附录A，附录B…。附录中的公式、图。附录中的公式、图和表的编号分别用和表的编号分别用A1，A2…系列；图A1，图A2…系列；表系列；表A1，表A2…系列。每个附录应有标题。每个附录应有标题。

\section{个人简历、在学期间发表的学术论文与研究成果}
个人简历包括出生年月日、获得学士、硕士学位的学校、时间等；学术论文研究成果按发表的时间顺序列出；研究成果可以是在学期间参加的研究项目、申请的专利或获奖情况等。

\section{致谢}
致谢是作者对该论文的形成作过贡献的组织或个人及参考文献的作者予以感谢的文字记载，语言要诚恳、恰当、简短。

